\documentclass[11pt]{article}

\usepackage{amsmath, amssymb}          % math (Lecture 3)
\usepackage[margin=1in]{geometry}      % margins
\usepackage{graphicx}                  % \includegraphics for figures
\usepackage{subcaption}                % side-by-side panels
\usepackage{booktabs}                  % professional tables (\toprule, etc.)
\usepackage{multirow}                  % cells spanning multiple rows
\usepackage{siunitx}                   % decimal-aligned numeric columns
\usepackage{algorithm}                 % the float (caption, label)
\usepackage{algpseudocode}             % the pseudocode inside it
\usepackage{natbib}                    % author-year citations
\usepackage{hyperref}                  % clickable refs and links (load late)
\usepackage[capitalise,noabbrev]{cleveref} % \cref writes "Figure 1" (load AFTER hyperref)
\hypersetup{colorlinks=true, linkcolor=blue, citecolor=blue, urlcolor=black}

\title{A Mini Paper: Figures, Tables, and Citations}
\author{Your Name}
\date{\today}

\begin{document}
\maketitle

\begin{abstract}
A short tour of the pieces that turn a LaTeX document into a paper:
cross-references, floats (figures and tables), algorithms, and a bibliography.
\end{abstract}

\noindent \textbf{Keywords:} LaTeX, reproducibility, typesetting

\section{Cross-references}\label{sec:refs}

Label things with \verb|\label{type:name}| and refer to them with \verb|\ref|
(or \verb|\eqref| for equations).
Never hard-code a number; LaTeX keeps them in sync, but you must compile
\emph{twice} for references to resolve.
This is Section~\ref{sec:refs}.
Consider the model
\begin{equation}\label{eq:slr}
  y_i = \beta_0 + \beta_1 x_i + \varepsilon_i,
  \qquad \varepsilon_i \overset{\text{iid}}{\sim} \mathcal{N}(0, \sigma^2),
\end{equation}
whose least-squares fit appears in \cref{fig:diagnostics}.
That reference used \verb|\cref|, which wrote the word ``Figure'' by itself.
We can also point back to Equation~\eqref{eq:slr} from anywhere.

\section{Tables}\label{sec:tables}

Inside \texttt{tabular}, \texttt{\&} separates columns and \verb|\\| ends a row;
the column spec sets alignment.
Prefer \texttt{booktabs} rules (\verb|\toprule|, \verb|\midrule|,
\verb|\bottomrule|) over vertical lines.
A table's caption goes \emph{above} it.

\begin{table}[ht!]
  \centering
  \caption{Simulation settings for the example in \cref{sec:tables}.}
  \label{tab:settings}
  \begin{tabular}{lrr}
    \toprule
    Scenario & Sample size $n$ & Noise SD $\sigma$ \\
    \midrule
    Baseline & 60  & 1.0 \\
    Larger   & 200 & 1.0 \\
    Noisier  & 60  & 2.5 \\
    \bottomrule
  \end{tabular}
\end{table}

\Cref{tab:estimates} is fancier: \verb|\multicolumn| spans a header across two
columns, \verb|\multirow| spans a label down two rows, and the \texttt{S} columns
from \texttt{siunitx} align on the decimal point (note the \texttt{\{braces\}}
around \texttt{S}-column headers).

\begin{table}[ht!]
  \centering
  \caption{Coefficient estimates, aligned on the decimal point.}
  \label{tab:estimates}
  \begin{tabular}{ll S[table-format=2.3] S[table-format=1.3]}
    \toprule
    & & \multicolumn{2}{c}{Estimates} \\
    \cmidrule(lr){3-4}
    Model & Prior & {Mean} & {SD} \\
    \midrule
    \multirow{2}{*}{OLS} & Flat  & 12.345 & 0.112 \\
                         & Ridge &  1.200 & 0.09  \\
    \midrule
    \multirow{2}{*}{GP}  & Flat  &  0.98  & 0.087 \\
                         & Ridge &  1.024 & 0.1   \\
    \bottomrule
  \end{tabular}
\end{table}

\verb|knitr::kable| and \texttt{modelsummary} in R, or \verb|df.to_latex()| in
pandas, give you the source code directly, which saves plenty of time.

\section{Figures}\label{sec:figures}

Load \texttt{graphicx} and use \verb|\includegraphics| inside a \texttt{figure}
float.
Size widths relative to \verb|\textwidth| so they stay portable, and refer to the
figure by number; floats move, so ``the figure below'' is a trap.
A figure's caption goes \emph{below} it.
The \texttt{subcaption} package gives individually-referenceable panels, so we can
cite \cref{fig:panel-b} on its own.

\begin{figure}[ht!]
  \centering
  \begin{subfigure}[b]{0.48\textwidth}
    \centering
    \includegraphics[width=\textwidth]{figures/scatter.pdf}
    \caption{Fitted model.}
    \label{fig:panel-a}
  \end{subfigure}
  \hfill
  \begin{subfigure}[b]{0.48\textwidth}
    \centering
    \includegraphics[width=\textwidth]{figures/residuals.pdf}
    \caption{Residual diagnostic.}
    \label{fig:panel-b}
  \end{subfigure}
  \caption{Simulated data with an ordinary least-squares fit, the model from
           \cref{eq:slr}. Generated by \texttt{make\_figure.py}.}
  \label{fig:diagnostics}
\end{figure}

\section{Algorithms}\label{sec:alg}

\texttt{algorithm} is the float; \texttt{algorithmic} is the pseudocode inside.
The \texttt{[1]} turns on line numbers, and \cref{alg:boot} is referenced like
anything else.

\begin{algorithm}[ht!]
  \caption{Nonparametric bootstrap \citep{efron1979bootstrap}}
  \label{alg:boot}
  \begin{algorithmic}[1]
    \Require data $x_1,\ldots,x_n$; replicates $B$
    \For{$b = 1$ \textbf{to} $B$}
      \State draw $x_1^*,\ldots,x_n^*$ with replacement from $x_1,\ldots,x_n$
      \State compute $\hat{\theta}^*_b$ from the resample
    \EndFor
    \State \Return the empirical distribution of $\hat{\theta}^*_1,\ldots,\hat{\theta}^*_B$
  \end{algorithmic}
\end{algorithm}

\section{Citations}\label{sec:cites}

Keep references in a \texttt{.bib} file and cite by key.
Use \verb|\citet| for a textual citation -- \citet{efron1979bootstrap} introduced
the bootstrap -- and \verb|\citep| for a parenthetical one
\citep{tibshirani1996regression}.
You can add notes inside the brackets, e.g.\ \citep[see][Ch.~7]{casella2002statistical},
name an author alone with \verb|\citeauthor| (\citeauthor{tibshirani1996regression}),
and cite software too \citep{wickham2016ggplot2}.
The list below is formatted by the chosen style; journals ship their own (drop in
their \texttt{.bst} and change \verb|\bibliographystyle|).

\appendix
\section{An Appendix}\label{app:extra}

\verb|\appendix| prints nothing; it just switches section numbering to letters.
This is \cref{app:extra}, and every reference still works.

\bibliographystyle{plainnat}
\bibliography{paper_example}

\end{document}
